\documentclass[aspectratio=169,10pt]{beamer}
\usetheme[progressbar=frametitle,numbering=fraction]{metropolis}
\usepackage[utf8]{inputenc}
\usepackage[T1]{fontenc}
\usepackage{lmodern}
\usepackage[spanish,es-noshorthands,es-tabla]{babel}
\usepackage{booktabs}
\usepackage{tikz}
\usetikzlibrary{positioning,arrows.meta}

\definecolor{guinda}{HTML}{7A1F2B}
\definecolor{acento}{HTML}{C8102E}
\setbeamercolor{frametitle}{bg=guinda}
\setbeamercolor{progress bar}{fg=acento}
\setbeamercolor{alerted text}{fg=acento}

% Generado por scripts/generar_metricas.py — no editar a mano
\newcommand{\NumPreguntas}{16}
\newcommand{\NumFragmentos}{1421}
\newcommand{\NumPaginas}{390}
\newcommand{\HitSim}{1{,}000}
\newcommand{\MrrSim}{1{,}000}
\newcommand{\PSim}{0{,}729}
\newcommand{\DivSim}{1{,}938}
\newcommand{\LatSim}{0{,}714}
\newcommand{\HitMMR}{1{,}000}
\newcommand{\MrrMMR}{1{,}000}
\newcommand{\PMMR}{0{,}625}
\newcommand{\DivMMR}{2{,}188}
\newcommand{\LatMMR}{0{,}612}
\newcommand{\PctCitas}{100\,\%}
\newcommand{\LatGen}{8{,}1}
\newcommand{\AnalisisRecuperacion}{Ambas estrategias recuperan al menos un fragmento relevante en todas las preguntas (Hit@6 de 1{,}000 con similitud y 1{,}000 con MMR). MMR eleva la diversidad de fuentes por consulta de 1{,}938 a 2{,}188, lo que aporta contexto comparado; ambas estrategias empatan en MRR. La menor P@6 de MMR es esperable: al penalizar la redundancia sustituye fragmentos contiguos del mismo artículo por fragmentos de otras fuentes.}
\newcommand{\AnalisisGeneracion}{En la inspección manual (\texttt{eval/resultados/respuestas.md}), las respuestas distinguen correctamente el carácter facultativo de la EIPD en el Perú frente a su obligatoriedad en el RGPD, y no se observaron artículos inventados. En 3 respuestas el modelo advierte de forma explícita que el corpus no cubre una parte de la pregunta (p. ej., que no existe una lista peruana de criterios de alto riesgo equivalente a WP248). Es el comportamiento esperado del prompt de fundamentación.}
\newcommand{\TablaCorpus}{LPDP & Ley N.º 29733 (Perú) & 17 & 63 \\
RLPDP & D.S. N.º 016-2024-JUS (Perú) & 60 & 200 \\
RGPD & Reglamento (UE) 2016/679 (Unión Europea) & 88 & 449 \\
WP248 & WP248 rev.01 (Unión Europea) & 22 & 79 \\
AEPD & AEPD (España) & 160 & 486 \\
NIST-PF & NIST Privacy Framework v1.0 (EE. UU.) & 43 & 144 \\}


\title{RAG normativo para la Evaluación de Impacto de Privacidad}
\subtitle{Avance TIF --- Agente de evaluación de impacto de privacidad}
\author{Dylan Antonio Zúñiga Huarca \and Jeferson Joao Basurco Casani \and Adrian Issac Mamani Quispe}
\institute{Auditoría de Sistemas · Escuela Profesional de Ingeniería de Sistemas · UNSA}
\date{\today}

\begin{document}
\maketitle

\begin{frame}{Problema}
\begin{columns}[T]
\column{.55\textwidth}
\begin{itemize}
  \item La \alert{EIPD / PIA} evalúa los riesgos de un tratamiento de datos \emph{antes} de ejecutarlo.
  \item El auditor debe cruzar la ley, el reglamento, el derecho comparado y las guías metodológicas: \alert{cientos de páginas}.
  \item Un LLM ``a secas'' \alert{alucina} artículos y plazos, algo inaceptable en una auditoría.
\end{itemize}
\column{.42\textwidth}
\begin{block}{Contexto peruano}
\small El nuevo Reglamento de la Ley 29733 (\textbf{D.S. 016-2024-JUS}) incorpora la EIPD en su \textbf{art. 40}, y su \textbf{art. 125} la considera atenuante de sanciones.
\end{block}
\end{columns}
\end{frame}

\begin{frame}{Objetivo del TIF y alcance de este avance}
\begin{center}
\begin{tikzpicture}[node distance=6mm,
  b/.style={draw,rounded corners,minimum height=9mm,text width=24mm,align=center,font=\small},
  on/.style={b,fill=acento!15,draw=acento,thick}, ar/.style={-Stealth,thick}]
\node[b] (a) {Descripción del tratamiento};
\node[b,right=of a] (b) {Tamizaje: ¿requiere EIPD?};
\node[b,right=of b] (c) {Riesgos\\prob. × impacto};
\node[b,right=of c] (d) {Medidas e informe EIPD};
\node[on,below=10mm of b, xshift=17mm, text width=40mm] (r) {\textbf{RAG normativo}\\(este avance)};
\draw[ar] (a)--(b); \draw[ar] (b)--(c); \draw[ar] (c)--(d);
\foreach \n in {b,c,d} \draw[ar,acento,dashed] (r) -- (\n);
\end{tikzpicture}
\end{center}
\small El agente completo se presenta en la unidad 3. Hoy entregamos su \alert{base de conocimiento}: responde con citas a la fuente, el artículo y la página.
\end{frame}

\begin{frame}{EIPD: Perú frente a la Unión Europea}
\small
\begin{tabular}{@{}p{.2\textwidth}p{.37\textwidth}p{.37\textwidth}@{}}
\toprule
 & \textbf{Perú (LPDP + RLPDP)} & \textbf{UE (RGPD)} \\
\midrule
Carácter & Facultativa (art. 40.1) & Obligatoria si hay alto riesgo (art. 35) \\
Supuestos & Datos sensibles, perfilamiento, personas vulnerables, grandes volúmenes & Perfilamiento con efectos jurídicos, datos sensibles a gran escala, vigilancia sistemática \\
Referencia & NTP-ISO/IEC 27005, NTP-ISO 31000 & WP248 (9 criterios), ISO/IEC 29134 \\
Incentivo & Atenuante de sanción (art. 125.4) & Consulta previa a la autoridad (art. 36) \\
\bottomrule
\end{tabular}
\end{frame}

\begin{frame}{¿Qué es RAG?}
\begin{columns}[c]
\column{.5\textwidth}
\textbf{R}etrieval-\textbf{A}ugmented \textbf{G}eneration (Lewis et al., 2020):
\begin{enumerate}
  \item \alert{Recuperar} los fragmentos relevantes de un corpus verificable.
  \item \alert{Aumentar} el prompt con esos fragmentos.
  \item \alert{Generar} una respuesta fundamentada y citada.
\end{enumerate}
\column{.45\textwidth}
\begin{block}{¿Por qué RAG y no fine-tuning?}
\small
\begin{itemize}
  \item Las normas cambian: basta con re-indexar.
  \item Trazabilidad: cada afirmación lleva su cita.
  \item Costo cero con el plan gratuito de Gemini.
\end{itemize}
\end{block}
\end{columns}
\end{frame}

\begin{frame}{Arquitectura}
\centering
\begin{tikzpicture}[node distance=5mm and 5mm,
  box/.style={draw,rounded corners=2pt,align=center,font=\footnotesize,minimum height=10mm,text width=24mm,fill=white},
  db/.style={box,fill=guinda!10}, llm/.style={box,fill=acento!12}, ar/.style={-Stealth,thick}]
\node[box] (pdf) {6 PDFs\\normativos};
\node[box,right=of pdf] (clean) {Extracción y limpieza (pypdf)};
\node[box,right=of clean] (chunk) {Fragmentos por artículo\\1200/200 car.};
\node[db,right=of chunk] (vec) {ChromaDB\\coseno · 768 d};
\node[box,below=12mm of pdf] (q) {Pregunta del auditor};
\node[llm,right=of q] (emb) {gemini-embedding-001};
\node[box,right=of emb] (mmr) {MMR\\k=6 de 24};
\node[llm,right=of mmr] (gen) {gemini-2.5-flash\\+ prompt citado};
\draw[ar] (pdf)--(clean); \draw[ar] (clean)--(chunk); \draw[ar] (chunk)--(vec);
\draw[ar] (q)--(emb); \draw[ar] (emb)--(mmr); \draw[ar] (vec)--(mmr); \draw[ar] (mmr)--(gen);
\node[font=\scriptsize,above=1mm of clean] {\textit{Ingesta (fuera de línea)}};
\node[font=\scriptsize,below=1mm of emb] {\textit{Consulta (en línea)}};
\end{tikzpicture}

\vspace{3mm}\small Python 3.12 · LangChain · ChromaDB · API de Gemini · Streamlit (demo)
\end{frame}

\begin{frame}{Corpus: \NumPaginas{} páginas, \NumFragmentos{} fragmentos}
\centering\small
\begin{tabular}{@{}llrr@{}}
\toprule
Sigla & Documento (jurisdicción) & Págs. & Frag. \\
\midrule
\TablaCorpus
\bottomrule
\end{tabular}
\end{frame}

\begin{frame}[fragile]{Fragmentación consciente de la norma}
\begin{columns}[T]
\column{.5\textwidth}
\begin{itemize}\small
  \item Los separadores priorizan \texttt{TÍTULO} → \texttt{CAPÍTULO} → \texttt{Artículo}.
  \item Se corrigen artefactos del PDF (``modifi catoria'', guiones de corte).
  \item Cada fragmento lleva sus metadatos: fuente, \alert{artículo}, \alert{página} y jurisdicción.
  \item Esto permite filtrar por jurisdicción (p. ej., solo Perú).
\end{itemize}
\column{.47\textwidth}
\begin{block}{Metadatos de un fragmento}
\scriptsize
\begin{verbatim}
{ "fuente": "RLPDP",
  "articulo": "40",
  "pagina": 24,
  "jurisdiccion": "Perú",
  "tipo": "norma" }
\end{verbatim}
\end{block}
\small Cita resultante: \alert{[RLPDP, art. 40, p. 24]}
\end{columns}
\end{frame}

\begin{frame}{Recuperación con MMR y prompt fundamentado}
\begin{columns}[T]
\column{.5\textwidth}
\textbf{Maximal Marginal Relevance}
{\small\[
d^{*}=\arg\max_{d_i}\Big[\lambda\,\mathrm{sim}(d_i,q)-(1-\lambda)\max_{d_j\in S}\mathrm{sim}(d_i,d_j)\Big]
\]}
\small Con $\lambda=0{,}6$ equilibra relevancia y diversidad: evita traer 6 trozos del mismo artículo.
\column{.47\textwidth}
\textbf{Reglas del prompt}
\begin{enumerate}\small
  \item Responder solo con los fragmentos.
  \item Citar cada afirmación con [n].
  \item Si falta información, decirlo.
  \item Separar la norma peruana vinculante del derecho comparado.
\end{enumerate}
\end{columns}
\end{frame}

\begin{frame}{Resultados de recuperación (\NumPreguntas{} preguntas, k = 6)}
\centering
\begin{tabular}{@{}lccccc@{}}
\toprule
Estrategia & Hit@6 & MRR & P@6 & Fuentes dist. & Lat. (s) \\
\midrule
Similitud & \HitSim & \MrrSim & \PSim & \DivSim & \LatSim \\
MMR & \HitMMR & \MrrMMR & \PMMR & \DivMMR & \LatMMR \\
\bottomrule
\end{tabular}

\vspace{5mm}
\begin{itemize}\small
  \item Generación: el \alert{\PctCitas} de las respuestas incluye citas; la latencia media es de \LatGen{} s.
  \item MMR diversifica las fuentes: aporta contexto comparado y cuesta algo de precisión.
  \item Primera iteración: Hit@6 = 0,938 y MRR = 0,818. Tras etiquetar los considerandos del RGPD y refinar la anotación se llega al techo, así que hace falta un conjunto más amplio y difícil.
\end{itemize}
\end{frame}

\begin{frame}[fragile]{Demostración: ¿es obligatoria la EIPD en el Perú?}
\small
\begin{block}{Respuesta del RAG}
En Perú, la EIPD es de carácter \textbf{facultativo} [2]. El titular del banco de datos \emph{puede} realizarla de manera previa al tratamiento, especialmente con datos sensibles, perfilamiento, personas vulnerables o grandes volúmenes [1]. Puede elaborarse tomando como referencia la NTP-ISO/IEC 27005 y la NTP-ISO 31000 [1].
\end{block}
\textbf{Fuentes:} [1] RLPDP, art. 40, p. 25 · [2] RLPDP, art. 40, p. 24 · \ldots

\vspace{2mm}
\normalsize\texttt{uv run streamlit run app.py} $\rightarrow$ demo en vivo
\end{frame}

\begin{frame}{Limitaciones y siguiente unidad}
\begin{columns}[T]
\column{.48\textwidth}
\textbf{Limitaciones}
\begin{itemize}\small
  \item Conjunto de evaluación pequeño: métricas indicativas.
  \item Tablas y columnas de los PDF pierden estructura.
  \item Las normas ISO/NTP no se incluyen por derechos de autor.
  \item Es un apoyo, no una asesoría legal.
\end{itemize}
\column{.48\textwidth}
\textbf{Unidad 3: agente completo}
\begin{itemize}\small
  \item Cuestionario del tratamiento.
  \item Tamizaje (WP248, art. 35 del RGPD, art. 40 del RLPDP).
  \item Matriz de riesgos (probabilidad × impacto).
  \item Informe EIPD automático con trazabilidad.
\end{itemize}
\end{columns}
\end{frame}

\begin{frame}[standout]
¿Preguntas?

\vspace{4mm}\small Código, informe IEEE y evaluación disponibles en el repositorio
\end{frame}

\end{document}
